\documentclass[11pt,a4paper]{article}
\usepackage{fontspec,amsmath,amssymb,amsthm,microtype}
\setmainfont{lmroman10-regular.otf}[BoldFont=lmroman10-bold.otf,ItalicFont=lmroman10-italic.otf,BoldItalicFont=lmroman10-bolditalic.otf]
\setmonofont{lmmono10-regular.otf}[Scale=MatchLowercase]
\usepackage[margin=25mm]{geometry}
\usepackage[unicode,hidelinks]{hyperref}
\setlength{\emergencystretch}{3em}
\setlength{\parskip}{0.4em}
\setlength{\parindent}{0pt}
\allowdisplaybreaks[2]
\raggedbottom
\title{\texorpdfstring{Normalize inverse mass in the actual kinetic equation}{Normalize inverse mass in the actual kinetic equation}}
\author{LangevinLean mathematical exposition}
\date{}
\begin{document}\maketitle
\noindent Source module: \texttt{Kinetic\allowbreak{}Langevin/\allowbreak{}Mass\allowbreak{}Scaling.\allowbreak{}lean}. The accompanying HTML pairs this mathematical exposition with the exact Lean statements and records the scope of the formalization.


The manuscript’s Lemma 1 removes inverse mass \(u>0\) by setting \(c=\sqrt u\), changing time to \(t/c\), and dividing velocity by \(c\). This proof works with the complete integral equation in an arbitrary complete normed space, so it applies to the constructed L\ensuremath{{}^{2}} Brownian solution without differentiating its noise.


\section{\texorpdfstring{Normalize the true phase drift}{Normalize the true phase drift}}


The continuous linear map \(D_a(v,x)=(av,x)\) changes velocity alone. Expanding both components proves \(c^{-1}D_{c^{-1}}R_{g,\gamma,c^2}(z)=R_{g,\gamma/c,1}(D_{c^{-1}}z)\). In the first component, the time and velocity factors cancel the original \(c^2\) force factor. In the second, the transformed position derivative is the normalized velocity. Lean verifies these as vector equalities using scalar action identities.



{\raggedright\noindent Lean declarations: \texttt{mass\allowbreak{}Phase\allowbreak{}Map}, \texttt{kinetic\_\allowbreak{}mass\_\allowbreak{}drift\_\allowbreak{}identity}.\par}

\section{\texorpdfstring{Transform the integral and its actual Brownian amplitude}{Transform the integral and its actual Brownian amplitude}}


Apply the bounded linear map to the solution identity at \(t/c\), pass it through the Bochner integral, and substitute \(r=s/c\). The transformed equation has friction \(\gamma/c\) and inverse mass one. The additive input becomes \(c^{-1}B_{t/c}\). For the original thermal amplitude \(\sqrt{2\gamma c^2}\), the coefficient identity gives \(c^{-1}\sqrt{2\gamma c^2}=\sqrt{2\gamma/c}\sqrt c\). The separately proved changed-clock Brownian process is \(\sqrt c W_{t/c}\), so the resulting equation has exactly the required fluctuation-dissipation coefficient. Positivity of \(c\) supplies both the interval inclusion and the square-root cancellations.



{\raggedright\noindent Lean declarations: \texttt{kinetic\_\allowbreak{}mass\_\allowbreak{}integral}, \texttt{kinetic\_\allowbreak{}mass\_\allowbreak{}noise\_\allowbreak{}amplitude}.\par}

\end{document}
